% Deklaracja członkowska Iławskiego Instytutu Idei (wersja dla kandydata, strona WWW).
% Pola formularza PDF pozwalają wypełnić deklarację na komputerze.
% Kompilacja: npm run deklaracja (XeLaTeX, dwa przebiegi).
% Marka: logo poziome i kolory z systemu marki Instytutu; kroje w fonts/ (SIL OFL 1.1).
\documentclass[11pt]{article}
\usepackage[a4paper,top=20mm,bottom=18mm,left=20mm,right=20mm,headheight=0pt,headsep=0pt,footskip=9mm]{geometry}
\usepackage{fontspec}
\usepackage{polyglossia}
\setdefaultlanguage{polish}
\usepackage{xcolor}
\usepackage{graphicx}
\usepackage{tikz}
\usepackage{eso-pic}
\usepackage{fancyhdr}
\usepackage{tabularx}
\usepackage{enumitem}
\usepackage{needspace}

% ---------- Kolory marki (tokens.css, jasny motyw) ----------
\definecolor{tekst}{HTML}{1D2320}
\definecolor{tekstwyciszony}{HTML}{5B625D}
\definecolor{linia}{HTML}{D9D7CC}
\definecolor{lab}{HTML}{2F7A29}
\definecolor{tlo}{HTML}{F7F6F1}
\definecolor{siatka}{HTML}{1F6B45}

% ---------- Kroje ----------
% sans: Noto Sans (krój systemowy), mono: IBM Plex Mono
\setmainfont{NotoSans}[Path=./fonts/, Extension=.ttf,
  UprightFont=*-Regular, BoldFont=*-SemiBold]
\setsansfont{NotoSans}[Path=./fonts/, Extension=.ttf,
  UprightFont=*-Regular, BoldFont=*-SemiBold]
\setmonofont{IBMPlexMono}[Path=./fonts/, Extension=.ttf,
  UprightFont=*-Regular, BoldFont=*-Bold]
\newfontfamily\naglowekfont{NotoSans-SemiBold.ttf}[Path=./fonts/]

\color{tekst}
\setlength{\parindent}{0pt}
\setlength{\parskip}{0pt}
\linespread{1.18}
\frenchspacing
\emergencystretch=3em
\pagestyle{empty}

% ---------- Motyw: papier milimetrowy w pasku nagłówka ----------
% Rysowany tylko na pierwszej stronie każdego formularza.
\newlength{\pasek}\setlength{\pasek}{42mm}
\newcommand{\siatkanaglowka}{%
  \AddToShipoutPictureBG*{%
  \begin{tikzpicture}[remember picture,overlay]
    \begin{scope}
      \clip (current page.north west) rectangle ([yshift=-\pasek]current page.north east);
      \draw[siatka!9!white, line width=0.25pt, step=5mm]
        (current page.north west) grid ([yshift=-\pasek]current page.north east);
    \end{scope}
    \draw[linia, line width=0.5pt]
      ([yshift=-\pasek]current page.north west) -- ([yshift=-\pasek]current page.north east);
  \end{tikzpicture}}}

% ---------- Stopka ----------
\newcommand{\stopkatekst}{}
\fancypagestyle{formularz}{%
  \fancyhf{}\renewcommand{\headrulewidth}{0pt}%
  \fancyfoot[L]{\ttfamily\footnotesize\color{tekstwyciszony}Iławski Instytut Idei}%
  \fancyfoot[R]{\ttfamily\footnotesize\color{tekstwyciszony}\stopkatekst}}
\pagestyle{formularz}

% ---------- Nagłówek formularza ----------
% Po lewej etykieta (mono), tytuł i podtytuł; po prawej logo poziome.
% #1 etykieta, #2 tytuł, #3 podtytuł, #4 tekst stopki
\newcommand{\naglowekformularza}[4]{%
  \clearpage\gdef\stopkatekst{#4}\siatkanaglowka
  \vspace*{-6mm}%
  \begin{minipage}[t]{0.66\linewidth}\raggedright
    \vspace{0pt}%
    \if\relax\detokenize{#1}\relax\else
      {\ttfamily\footnotesize\color{lab}\MakeUppercase{#1}\par}
      \vspace{2mm}%
    \fi
    {\naglowekfont\fontsize{19}{23}\selectfont #2\par}
    \vspace{1mm}%
    {\fontsize{11.5}{14}\selectfont\color{tekstwyciszony}#3\par}
  \end{minipage}\hfill
  \begin{minipage}[t]{0.3\linewidth}\raggedleft
    \vspace{0pt}%
    \includegraphics[height=13mm]{iii-logo-poziome-kolor.pdf}
  \end{minipage}\par
  \vspace{11mm}}

% ---------- Sekcja ----------
\newcommand{\sekcja}[1]{%
  \needspace{15mm}\vspace{3mm}%
  {\naglowekfont\fontsize{12.5}{15}\selectfont #1\par}
  \vspace{1.5mm}}

% ---------- Pole do wypełnienia ----------
% Etykieta nad linią, miejsce na pismo odręczne.
% Pole formularza PDF nad linią: \nazwa{id} przed \pole włącza pole.
\newcommand{\polenazwa}{}
\newcommand{\nazwa}[1]{\gdef\polenazwa{#1}}
\newcommand{\pole}[2][\linewidth]{%
  \begin{minipage}[b]{#1}
    {\small\raggedright\color{tekstwyciszony}#2\par}
    \vspace{5.5mm}%
    \ifx\polenazwa\empty\else
      \makebox[0pt][l]{\smash{\raisebox{0.4mm}{\TextField[name=\polenazwa,
        width=\linewidth, height=5.5mm, charsize=10pt, bordercolor={},
        backgroundcolor={}, color=0.114 0.137 0.125]{}}}}%
      \gdef\polenazwa{}%
    \fi
    {\color{tekst}\rule{\linewidth}{0.4pt}}
  \end{minipage}}
\newcommand{\polewiersz}[1]{\par\vspace{1.5mm}\pole{#1}\par\vspace{3.5mm}}

% Dwa pola obok siebie
\newcommand{\dwapola}[2]{%
  \par\vspace{1.5mm}\pole[0.48\linewidth]{#1}\hfill\pole[0.48\linewidth]{#2}\par\vspace{3.5mm}}

% Wpis w wierszu tekstu (liczby, numery, daty)
\newcommand{\wpis}[1]{\,\rule[-0.8mm]{#1}{0.4pt}}
\newcommand{\wiersz}[1]{\par\vspace{2.5mm}%
  {\fontsize{10.25}{14.5}\selectfont #1\par}\vspace{2.5mm}}

% Wynik głosowania w jednym wierszu
\newcommand{\glosowanie}{\wiersz{Za:\wpis{24mm}\hfill
  Przeciw:\wpis{24mm}\hfill Wstrzymujące się:\wpis{24mm}}}

% Liczby obecności
\newcommand{\obecnosc}[1]{%
  \wiersz{Liczba członków uprawnionych przed przyjęciem nowej osoby:\hfill\wpis{30mm}}
  \vspace{-1mm}\wiersz{Liczba obecnych uprawnionych do głosowania:\hfill\wpis{30mm}}}

% Tekst zwykły
\newcommand{\tresc}[1]{{\fontsize{10.25}{14}\selectfont #1\par}\vspace{2.6mm}}

% Uwaga małym tekstem
\newcommand{\uwaga}[1]{{\fontsize{9}{12}\selectfont\color{tekstwyciszony}#1\par}}
\newcommand{\info}[1]{{\fontsize{8.5}{11}\selectfont\raggedright\color{tekstwyciszony}#1\par}}

\usepackage[pdftitle={Deklaracja członkowska — Iławski Instytut Idei},
  pdfauthor={Stowarzyszenie zwykłe „IŁAWSKI INSTYTUT IDEI”},
  pdfsubject={Deklaracja członkowska},
  pdflang={pl-PL},hidelinks]{hyperref}

\begin{document}
\begin{Form}

\naglowekformularza{}
  {Deklaracja członkowska}
  {Stowarzyszenie zwykłe „IŁAWSKI INSTYTUT IDEI”}
  {instytut.ilawa.pl/regulamin}

\tresc{Zwracam się do Zebrania Członków o przyjęcie mnie w poczet członków
Stowarzyszenia zwykłego „IŁAWSKI INSTYTUT IDEI”.}

\nazwa{imie_nazwisko}\polewiersz{Imię i nazwisko}
\par\vspace{1.5mm}\nazwa{kontakt}\pole[0.62\linewidth]{Kontakt do korespondencji (adres lub e-mail)}\hfill
\nazwa{telefon}\pole[0.34\linewidth]{Telefon (dobrowolnie)}\par\vspace{3.5mm}

\vspace{1mm}
\tresc{Oświadczam, że posiadam pełną zdolność do czynności prawnych i~nie
jestem pozbawiony/\allowbreak pozbawiona praw publicznych.
Zapoznałem/\allowbreak zapoznałam się z~Regulaminem Stowarzyszenia
i~zobowiązuję się do jego przestrzegania, wykonywania uchwał władz
Stowarzyszenia oraz wypełniania obowiązków członkowskich, w~tym aktywnego
udziału w~realizacji celów Stowarzyszenia i~opłacania ustalonych składek.}

\par\vspace{1.5mm}\nazwa{miejscowosc_data}\pole[0.28\linewidth]{Miejscowość i data}\hfill
\nazwa{skladka}\pole[0.3\linewidth]{Deklarowana wysokość składki członkowskiej}\hfill
\pole[0.36\linewidth]{Podpis kandydata/\allowbreak kandydatki}\par\vspace{3.5mm}

\sekcja{Rekomendacja dwóch członków Stowarzyszenia}
\tresc{Rekomendujemy przyjęcie wyżej wskazanej osoby do Stowarzyszenia.}

\vspace{1mm}{\ttfamily\footnotesize\color{lab}REKOMENDACJA 1\par}
\par\vspace{1.5mm}\nazwa{rekomendujacy_1}\pole[0.48\linewidth]{Imię i nazwisko członka rekomendującego}\hfill
\pole[0.48\linewidth]{Data i podpis}\par\vspace{3.5mm}

\vspace{1mm}{\ttfamily\footnotesize\color{lab}REKOMENDACJA 2\par}
\par\vspace{1.5mm}\nazwa{rekomendujacy_2}\pole[0.48\linewidth]{Imię i nazwisko członka rekomendującego}\hfill
\pole[0.48\linewidth]{Data i podpis}\par\vspace{3.5mm}

\vfill
{\color{linia}\rule{\linewidth}{0.5pt}}\par\vspace{2.5mm}
\begin{minipage}[t]{0.47\linewidth}
  {\ttfamily\footnotesize\color{lab}JAK TO DZIAŁA\par}\vspace{1mm}
  \info{Wypełnij deklarację i~poproś dwóch członków Stowarzyszenia
  o~rekomendację. Zebranie Członków rozpatruje deklarację uchwałą
  w~ciągu 30 dni od jej złożenia (pkt~10 Regulaminu). Wysokość
  składki deklarujesz w~porozumieniu z~władzami Stowarzyszenia.
  Sugerowana składka pokrywa bieżące koszty przestrzeni.\par\vspace{1.5mm}
  Regulamin: instytut.ilawa.pl/regulamin\par
  Pytania: kontakt@instytut.ilawa.pl}
\end{minipage}\hfill
\begin{minipage}[t]{0.47\linewidth}
  {\ttfamily\footnotesize\color{lab}TWOJE DANE\par}\vspace{1mm}
  \info{Administratorem danych jest Stowarzyszenie zwykłe „IŁAWSKI
  INSTYTUT IDEI” z~siedzibą w~Iławie, kontakt@instytut.ilawa.pl.
  Dane z~deklaracji przetwarzamy, aby rozpatrzyć wniosek i~prowadzić
  listę członków. Przechowujemy je przez czas członkostwa i~okres
  wymagany przepisami.
  Masz prawo dostępu do danych, ich sprostowania, usunięcia,
  ograniczenia przetwarzania, sprzeciwu i~skargi do Prezesa UODO.}
\end{minipage}

\end{Form}
\end{document}
